% !TEX root = ../../hw01.tex
% Homework 01, Exercise 7. Shared solution.

\begin{proof}[Solution]
Suppose that
\[
  r_1\tau+\cdots+r_n\tau^n=0,
  \qquad r_1,\ldots,r_n\in\Q.
\]
We want to show that every coefficient is zero. Suppose, for a
contradiction, that at least one $r_j\ne0$.

Write $r_j=a_j/b_j$, where $a_j\in\Z$ and $b_j\in\N$.
Put $D=b_1\cdots b_n>0$. We multiply through by $D$ to clear
all the denominators and obtain
\[
  (Dr_1)\tau+\cdots+(Dr_n)\tau^n=0.
\]
Each $Dr_j$ is an integer, and at least one is nonzero since
$D\ne0$. Therefore
\[
  P(X)=(Dr_1)X+\cdots+(Dr_n)X^n
\]
is a nonzero polynomial with integer coefficients satisfying
$P(\tau)=0$. This would make $\tau$ algebraic, contradicting
the assumption that it is transcendental.

Thus $r_1=\cdots=r_n=0$. Since $n$ was arbitrary, the numbers
$\tau,\ldots,\tau^n$ are linearly independent over $\Q$ for every
$n\in\N$.
\marginnote{Including a constant term does not change the argument.
  Thus $1,\tau,\ldots,\tau^n$ are also linearly independent over $\Q$.}
\end{proof}
